% !TEX TS-program = xelatex
% !TEX encoding = UTF-8 Unicode
% !Mode:: "TeX:UTF-8"

\documentclass{resume}
\usepackage{zh_CN-Adobefonts_external} % Simplified Chinese Support using external fonts (./fonts/zh_CN-Adobe/)
% \usepackage{NotoSansSC_external}
% \usepackage{NotoSerifCJKsc_external}
% \usepackage{zh_CN-Adobefonts_internal} % Simplified Chinese Support using system fonts
\usepackage{linespacing_fix} % disable extra space before next section
\usepackage{cite}

\begin{document}
\pagenumbering{gobble} % suppress displaying page number

\name{吴童}

\basicInfo{
  \email{wutong1109@stu.pku.edu.cn} \textperiodcentered\ 
  \phone{(+86) 159-5911-5743} \textperiodcentered\ 
  \linkedin[billryan8]{https://www.linkedin.com/in/billryan8}}
 
\section{\faGraduationCap\  教育背景}
\datedsubsection{\textbf{北京大学}, 北京}{2022 -- 至今}
\textit{本科生}\ 计算机科学与技术, 预计 2026 年 6 月毕业
\datedsubsection{\textbf{福州第一中学}, 福建, 福州}{2019 -- 2022}
\textit{高中生}

\section{\faUsers\ 实习/项目经历}
\datedsubsection{\textbf{梁云老师课题组} 北京大学}{2024年3月 -- 2024年5月}
\role{科研轮转}{指导老师: 梁云教授教授}
大模型推理与KVCache压缩
\begin{itemize}
  \item 学习LLM推理、KVCache基本知识和已有的压缩算法
  \item 以第四作者身份发表论文并被NIPS'24接收，首次提出可召回的压缩算法
\end{itemize}

\datedsubsection{\textbf{PKU-DAIR} 北京大学}{2024年8月 -- 2025年1月}
\role{本科生实习}{指导老师: 崔斌教授}
AI系统、分布式与并行训练
\begin{itemize}
  \item 参加VLDB'24会议
  \item 学习分布式并行训练相关知识与论文
\end{itemize}

\section{\faCogs\ 编程技能}
% increase linespacing [parsep=0.5ex]
\begin{itemize}[parsep=0.5ex]
  \item 语言: C, C++, Python, CUDA
  \item 平台: Linux
  \item 深度学习库：PyTorch
\end{itemize}

\section{\faHeartO\ 学术水平}
\datedline{\textit{入选2022级计算机拔尖班}}{2023.9}
\datedline{\textit{绩点}：3.80，top10%}{截至2024.12}
\datedline{\textit{学习优秀奖}}{2023.9}
\datedline{\textit{三好学生}}{2024.9}
\datedline{\textit{永旺奖学金}}{2024.9}
\datedline{\textit{论文发表}：[ArkVale: Efficient Generative LLM Inference with Recallable Key-Value Eviction](https://neurips.cc/virtual/2024/poster/96635)}


\section{\faInfo\ 其他}
% increase linespacing [parsep=0.5ex]
\begin{itemize}[parsep=0.5ex]
  \item GitHub: https://github.com/Rachmanino
  \item 语言: 英语 - 熟练(TOEFL 109)
\end{itemize}

%% Reference
%\newpage
%\bibliographystyle{IEEETran}
%\bibliography{mycite}
\end{document}
